\section{Conclusiones y Trabajo Futuro}

La revisión de la literatura  deja algo claro y es que para que una aplicación con GenAI o RAG sea sostenible en el largo plazo no solo importa la exactitud del modelo de IA sino que importa tanto o más la calidad del código, aquí el problema radica en que aunque para el software tradicional se cuenta con herramientas maduras para medir la deuda técnica, el ecosistema de la IA generativa aun carece de estos estándares y marcos de evaluación para auditar las arquitecturas.

Hoy en día los errores de integración en los sistemas RAG se buscan a mano y de forma reactiva, lo cual va generar un mayor consumo de recursos tanto computacionales como humanos, hasta el momento no se encuentra el desarrollo de un analizador estático que opere de forma proactiva para auditar la arquitectura especifica de este tipo de código de forma preventiva.

Para abordar esto, y en alineación con los objetivos de esta investigación, el trabajo futuro se centrará en el desarrollo de una herramienta automatizada (Linter/CLI) ejecutable en Python. Este instrumento utilizará el análisis de Árboles de Sintaxis Abstracta (AST) para escanear el código fuente de aplicaciones GenAI/RAG y detectar, en tiempo de diseño, violaciones estructurales y antipatrones basados en taxonomías recientes de la literatura, como la propuesta por Shao et al. \cite{shao_are_2024}

Finalmente, para garantizar la validez científica de la herramienta, se proyecta evaluar su precisión frente a conjuntos de datos de defectos empíricamente comprobados, como la biblioteca HYDRANGEA (biblioteca construida por el propio equipo de Shao et al.) \cite{shao_are_2024}. A partir de los hallazgos automatizados, se buscará establecer una evaluación estadística que correlacione la presencia de estos antipatrones con el deterioro de métricas tradicionales de diseño, tales como el Índice de Mantenibilidad (Maintainability Index), el Acoplamiento entre Objetos (CBO) y la Complejidad Ciclomática \cite{chidamber_metrics_1994}.